\documentclass[a4paper]{article}
\usepackage{amsmath}
\usepackage{amssymb}

\begin{document}

\large Auteur: Sadia Nazary \\
\large Studierichting: Informatica \\
\large Oefeningenreeks 3A nr. 12.2 \\

\medskip

\normalsize

Gegeven \\

$f(x)=(x^2+5x-6)(x^2-2x+4)$ \\

Methode 1 eerst uitwerken dan afleiden \\

$f(x)=(x^2+5x-6)(x^2-2x+4)$ \\

$=x^2(x^2-2x+4)+5x(x^2-2x+4)-6(x^2-2x+4)$ \\

$=x^4-2x^3+4x^2+5x^3-10x^2+20x-6x^2+12x-24$ \\

$=x^4+3x^3-12x^2+32x-24$ \\

Afgeleide nemen \\

$f'(x)=4x^3+9x^2-24x+32$ \\

Methode 2 productregel \\

$u=x^2+5x-6 \quad v=x^2-2x+4$ \\

$f'(x)=u'v+uv'$ \\

$u'=2x+5 \quad v'=2x-2$ \\

$f'(x)=(2x+5)(x^2-2x+4)+(x^2+5x-6)(2x-2)$ \\

$=2x^3-4x^2+8x+5x^2-10x+20$ \\

$+2x^3+10x^2-12x-2x^2-10x+12$ \\

$=4x^3+9x^2-24x+32$ \\

Beide methodes geven dezelfde afgeleide \\

$f'(x)=4x^3+9x^2-24x+32$ \\

\begin{thebibliography}{999}
\end{thebibliography}

\end{document}